\chapter{Arquitetura de Software}
\label{cap:arquitetura-teoria}

\begin{edital}
Arquitetura de software; interoperabilidade; arquitetura e linguagem
orientada a serviços; web services; mensageria; API, Swagger; arquitetura
orientada a objetos; aplicações para ambiente web; servidor de aplicações x
servidor web; internet/extranet/intranet/portal; arquitetura hexagonal;
microsserviços (orquestração e API gateway); containers; transações
distribuídas.
\end{edital}

\section{Servidor web x servidor de aplicação}

\begin{conceito}[Onde a requisição HTTP para, e onde a lógica de negócio roda]
\begin{center}
\begin{tabular}{@{}p{2.4cm}p{4.6cm}p{4.6cm}@{}}
\toprule
& \textbf{Servidor web} & \textbf{Servidor de aplicação} \\
\midrule
Função & atende HTTP, serve conteúdo estático/dinâmico & executa
\textbf{lógica de negócio}, integra com banco de dados \\
Exemplos & Apache, NGINX & JBoss/WildFly, WebLogic \\
\bottomrule
\end{tabular}
\end{center}

No mundo Java a distinção é literal, não só uma divisão de responsabilidade
conceitual: o servidor de aplicação implementa a especificação
\textbf{Java EE/Jakarta EE completa} --- ele traz o \textbf{container web}
(que processa Servlet/JSP, a mesma coisa que um servidor web faz)
\textbf{e, além disso}, o \textbf{container EJB} para os componentes de
negócio, mais transações distribuídas, JMS (mensageria) e injeção de
dependência. O Tomcat, por causa disso, é um caso interessante: ele é
servidor web \textbf{mais} container web, mas \textbf{não} é um servidor de
aplicação completo, porque não tem o container EJB.
\end{conceito}

\begin{cuidado}
Servidor web \textbf{não} processa lógica de negócio --- ele recebe a
requisição HTTP e serve conteúdo; quem executa a regra de negócio é o
servidor de aplicação. Servidor de aplicação não é “sempre mais leve”
(é o oposto: carrega containers a mais), não “dispensa a JVM” (roda sobre
ela) e não “deixa de suportar HTTP” (suporta, porque contém o próprio
container web dentro de si).
\end{cuidado}

\section{Estilos de arquitetura: uma sequência histórica}

\begin{conceito}[Cada estilo nasceu resolvendo a dor do anterior]
Os estilos de arquitetura não são uma lista de opções paralelas e
equivalentes --- eles formam uma \textbf{sequência histórica}, e entender
por que cada um surgiu é o que permite responder a uma questão de cenário
(“qual estilo resolve esta dor”), em vez de decorar definições soltas.

O \textbf{monólito} incomodava porque uma equipe travava a outra: subir a
correção de uma única tela obrigava a reimplantar o sistema inteiro. A
\textbf{SOA} respondeu quebrando o sistema em \textbf{serviços} com
contrato explícito e baixo acoplamento --- mas concentrou a inteligência de
integração num \textbf{barramento} central (o ESB), e esse barramento
virou o novo gargalo e ponto único de falha. Os \textbf{microsserviços}
responderam à SOA invertendo a filosofia: em vez de um barramento
inteligente com canais burros, a ideia virou \textbf{canais burros e
serviços espertos} --- cada serviço decide sozinho, sem depender de um
intermediário central pesado.

\begin{itemize}
\item \textbf{Monolítica}: tudo num artefato único, implantado junto.
Simples de começar; acopla e dificulta escalar partes isoladamente.
\item \textbf{Cliente-servidor, N camadas, P2P, barramento de mensagens}:
formas intermediárias de distribuir responsabilidade entre nós.
\item \textbf{SOA (orientada a serviços)}: serviços reutilizáveis, contrato
explícito, baixo acoplamento, interoperabilidade entre sistemas
heterogêneos. Web services são a tecnologia comum de implementação. O
barramento típico é o \textbf{ESB (Enterprise Service Bus)}: o
intermediário que \textbf{roteia}, \textbf{transforma} (converte
formato/protocolo entre sistemas que não se falam nativamente) e
\textbf{orquestra} as mensagens entre serviços.
\item \textbf{Microsserviços}: serviços pequenos, autônomos,
\textbf{implantáveis independentemente}, cada um com \textbf{seu próprio
banco de dados} (não compartilham base --- isso é o que reduz o
acoplamento). Comunicação leve, tipicamente REST ou mensageria.
\item \textbf{Arquitetura hexagonal (\textit{Ports \& Adapters})}: separa a
\textbf{lógica de negócio} do mundo externo através de \textbf{portas}
(interfaces que o núcleo define) e \textbf{adaptadores} (implementações
concretas dessas portas) --- permite trocar UI, banco de dados ou serviço
externo sem tocar no núcleo de negócio.
\end{itemize}
\end{conceito}

\begin{conceito}[O preço dos microsserviços --- por que nem sempre valem a pena]
Todo ganho de microsserviço vem com uma conta correspondente:

\begin{center}
\begin{tabular}{@{}p{5.4cm}p{7cm}@{}}
\toprule
\textbf{Ganho} & \textbf{Preço correspondente} \\
\midrule
implantar um serviço sem tocar nos outros & precisa de automação de
\textit{deploy} e observabilidade distribuída --- sem isso, achar a causa
de um erro fica pior do que num monólito \\
escalar só a parte que aperta & a chamada entre serviços virou
\textbf{chamada de rede}: pode falhar, tem latência, exige repetição e
\textit{timeout} \\
cada serviço com seu próprio banco & acabou o \texttt{JOIN} entre eles, e
acabou a transação ACID atravessando os dois --- daí a necessidade de
padrões como Saga e a aceitação de consistência eventual (ver
Seção~\ref{sec:transacoes-dist}) \\
equipes autônomas por serviço & só compensa se as equipes \textbf{forem}
realmente várias; para um time pequeno, é custo sem benefício
correspondente \\
\bottomrule
\end{tabular}
\end{center}

A regra prática: microsserviço resolve um problema \textbf{organizacional}
(várias equipes que precisam trabalhar em paralelo sem se travar) e de
\textbf{escala independente} (um componente específico recebe muito mais
carga que os outros). Se o cenário descreve sistema pequeno, equipe única
ou forte necessidade de consistência transacional entre módulos, um
monólito (eventualmente modular internamente) costuma ser a resposta mais
apropriada.
\end{conceito}

\begin{cuidado}
Microsserviços \textbf{não} compartilham o mesmo banco de dados ---
compartilhar aumentaria o acoplamento, contrariando o princípio central do
estilo. Baixo acoplamento significa que uma mudança em um serviço
\textbf{não} obriga mudança em outro. Um item que descreve “o barramento
que intermedeia, roteia e transforma mensagens entre serviços” está
descrevendo um \textbf{ESB} --- não confundir com ETL (move dados para um
\textit{data warehouse}), CDN (entrega conteúdo estático) ou DNS (resolve
nomes).
\end{cuidado}

\subsection{Camadas lógicas (\textit{layers}) x camadas físicas (\textit{tiers})}

\begin{conceito}[Duas ideias que o português funde numa palavra só]
\begin{center}
\begin{tabular}{@{}p{3.4cm}p{8.6cm}@{}}
\toprule
\textbf{Layer (lógica)} & como as \textbf{responsabilidades} do software
estão organizadas: apresentação, negócio, persistência \\
\textbf{Tier (física)} & em \textbf{quantos nós} o software está implantado:
máquinas, processos, servidores diferentes \\
\bottomrule
\end{tabular}
\end{center}

Uma coisa não obriga a outra. Uma aplicação pode ter \textbf{três layers e
um único tier}: as três responsabilidades bem separadas no código-fonte
(classes de apresentação, de negócio, de persistência, cada uma no seu
lugar), mas tudo rodando dentro do mesmo processo, na mesma máquina --- e a
aplicação continua sendo, do ponto de vista de implantação,
\textbf{monolítica}. O inverso também existe: um cliente-servidor com dois
tiers físicos pode manter apresentação e negócio bem separados
logicamente dentro de cada tier.
\end{conceito}

\begin{cuidado}
Não inverta as definições (layer é organização lógica de responsabilidade;
tier é distribuição física em nós). E desconfie de “três camadas lógicas
são \textbf{necessariamente} implantadas em três nós físicos” --- não são;
o MVC, por exemplo, é um padrão de organização \textbf{lógica}
(apresentação/modelo/controle), não um padrão de implantação em tiers.
\end{cuidado}

\section{Integração: REST, SOAP e web services}

\begin{conceito}[O que REST é de verdade --- recurso, representação, verbo]
REST não é “API que devolve JSON” --- é um \textbf{estilo arquitetural}
com uma ideia central: tudo que interessa é um \textbf{recurso},
identificado por uma URI, e o cliente nunca manipula o recurso diretamente
--- ele troca \textbf{representações} dele (JSON, XML, HTML). Dessa ideia
central decorrem as restrições do estilo:

\begin{itemize}
\item \textbf{Interface uniforme}: o conjunto de verbos é fixo e igual para
todo recurso --- \texttt{GET} lê, \texttt{POST} cria, \texttt{PUT}
substitui por inteiro, \texttt{PATCH} altera em parte, \texttt{DELETE}
remove. Por isso a URI deve nomear uma \textbf{coisa}, não uma
\textbf{ação}: \texttt{/pedidos/42} está de acordo com o estilo;
\texttt{/buscarPedido?id=42} contraria, porque põe o verbo na URI quando
ele já deveria estar só no método HTTP.
\item \textbf{Sem estado (\textit{stateless})}: cada requisição carrega
tudo o que é necessário para ser entendida sozinha, sem depender de uma
sessão guardada no servidor. O ganho não é ideológico: se o servidor não
guarda sessão, \textbf{qualquer} instância do servidor consegue atender
\textbf{qualquer} requisição --- e é exatamente isso que permite colocar um
balanceador de carga na frente e escalar horizontalmente sem precisar de
“afinidade de sessão” (a mesma requisição sempre indo para a mesma
instância).
\item \textbf{Cliente-servidor, cache, sistema em camadas}: o cliente não
precisa saber se está falando com o servidor final ou com um
intermediário --- por isso um proxy, um gateway ou uma CDN cabem no
caminho da requisição sem quebrar nada do lado do cliente.
\end{itemize}

A corrente de raciocínio que sustenta “por que REST escala bem”: ausência
de estado no servidor $\to$ qualquer instância pode atender qualquer
requisição $\to$ escala horizontal sem restrição de afinidade.
\end{conceito}

\begin{conceito}[REST x SOAP]
\begin{center}
\begin{tabular}{@{}p{2.6cm}p{4.6cm}p{4.6cm}@{}}
\toprule
& \textbf{REST} & \textbf{SOAP} \\
\midrule
Natureza & estilo arquitetural, sobre HTTP & protocolo baseado em XML \\
Formato & JSON (comum), leve & XML (envelope), verboso \\
Contrato & OpenAPI/Swagger & WSDL \\
Estado & \textbf{stateless} por definição & pode incorporar padrões WS-* \\
Vantagem & leve, escala, independe de plataforma & contratos formais,
padrões de segurança (WS-Security) \\
\bottomrule
\end{tabular}
\end{center}
\end{conceito}

\begin{cuidado}
\textbf{Cacheável} é uma restrição \textbf{distinta} de \textit{stateless}
--- a resposta se declara cacheável ou não, e quem a guarda pode ser o
\textbf{cliente} ou um \textbf{intermediário} (proxy, gateway, CDN). Cache
no caminho é \textbf{permitido} pelo REST, é o que deixa uma CDN servir
conteúdo estático de um nó próximo do usuário. O que \textbf{não} pode
ficar guardado no servidor é o \textbf{estado da sessão}, não o cache da
resposta --- não confunda os dois. \textbf{PUT} substitui o recurso
\textbf{inteiro}; \textbf{PATCH} altera só uma \textbf{parte} dele. Dizer
que “SOAP não tem contrato formal” contradiz o próprio SOAP, que usa WSDL
como contrato.
\end{cuidado}

\begin{conceito}[UDDI, WSDL, SOAP: três papéis, não um sinônimo do outro]
O trio clássico dos web services WS-*: \textbf{SOAP} é o formato de
\textbf{mensagem} trocada; \textbf{WSDL} é o \textbf{contrato} (descreve
o que o serviço oferece, quais operações, quais parâmetros); \textbf{UDDI}
é o diretório (legado) para \textbf{descoberta/registro} de web services
--- onde procurar quem oferece o quê. \textbf{Swagger/OpenAPI} cumpre, no
mundo REST, o papel que o WSDL cumpre no mundo SOAP: \textbf{OpenAPI} é a
especificação padrão para documentar/contratar APIs REST (endpoints,
parâmetros, respostas); \textbf{Swagger} é o conjunto de ferramentas
(Swagger UI, editor, geração de código) construído em cima dela.
\end{conceito}

\subsection{Mensageria: comunicação assíncrona}

\begin{conceito}[Desacopla quem envia de quem processa]
Numa chamada REST síncrona, quem envia \textbf{espera} a resposta. Na
mensageria, quem produz a mensagem \textbf{não espera} quem a consome ---
os dois lados só compartilham um intermediário, o \textbf{broker}
(RabbitMQ, Kafka).

\begin{center}
\begin{tabular}{@{}p{3.2cm}p{4.6cm}p{3cm}@{}}
\toprule
\textbf{Modelo} & \textbf{Como funciona} & \textbf{Exemplo} \\
\midrule
Fila (\textit{queue}) & mensagem consumida por \textbf{um} consumidor
(\textit{point-to-point}) & RabbitMQ, SQS \\
\textit{Publish/Subscribe} (tópico) & mensagem entregue a \textbf{vários}
assinantes ao mesmo tempo & Kafka, SNS \\
\bottomrule
\end{tabular}
\end{center}

\textbf{Apache Kafka} é uma plataforma de \textbf{\textit{streaming} de
eventos}: produtores publicam em \textbf{tópicos} (particionados para
paralelismo), consumidores leem organizados em \textbf{grupos}, e a
plataforma guarda o \textbf{log de eventos} (o que permite reprocessar
mensagens antigas, diferente de uma fila tradicional que descarta a
mensagem ao ser consumida).

Benefícios da mensageria: \textbf{desacoplamento} (produtor e consumidor
não precisam estar disponíveis ao mesmo tempo), \textbf{absorção de picos}
(a fila funciona como \textit{buffer} quando o consumidor não acompanha o
ritmo do produtor), \textbf{resiliência} (se o consumidor cair, a mensagem
espera na fila). Casa naturalmente com \textbf{microsserviços} e com os
padrões \textbf{Saga} (Seção~\ref{sec:transacoes-dist}) e
\textbf{Event Sourcing} (guardar o histórico de eventos como fonte da
verdade, em vez de só o estado final).
\end{conceito}

\begin{cuidado}
Fila entrega para \textbf{um} consumidor; tópico/\textit{pub-sub} entrega
para \textbf{vários} --- não inverta os dois. Mensageria é sempre
\textbf{assíncrona}, ao contrário de uma chamada REST tradicional, que é
síncrona.
\end{cuidado}

\section{Ambientes de rede corporativa}

\begin{conceito}
\begin{center}
\begin{tabular}{@{}p{2.4cm}p{9cm}@{}}
\toprule
\textbf{Termo} & \textbf{Alcance} \\
\midrule
Internet & rede pública global \\
Intranet & rede interna da organização (acesso só de funcionários) \\
Extranet & acesso controlado estendido a parceiros/clientes externos
autorizados --- o meio-termo entre intranet e internet \\
Portal & ponto único que centraliza acesso a informação/serviços diversos \\
\bottomrule
\end{tabular}
\end{center}
\end{conceito}

\section{Nuvem e escalabilidade}
\label{sec:nuvem-teoria}

\begin{conceito}[IaaS, PaaS, SaaS --- onde passa a linha de responsabilidade]
Os três modelos respondem à mesma pergunta: \textbf{até onde vai o
provedor, e onde começa a sua responsabilidade}. A pilha completa de uma
aplicação é sempre a mesma --- rede, armazenamento, servidores,
virtualização, sistema operacional, \textit{middleware}, tempo de
execução, dados, aplicação --- e cada modelo apenas move a linha de corte
mais para cima na pilha.

\begin{center}
\begin{tabular}{@{}p{2cm}p{4.6cm}p{5.4cm}@{}}
\toprule
& \textbf{Você cuida de} & \textbf{Exemplo típico} \\
\midrule
\textbf{IaaS} & sistema operacional para cima: \textit{patch}, tempo de
execução, aplicação, dados & máquina virtual alugada \\
\textbf{PaaS} & aplicação e dados; o resto é responsabilidade do provedor
& plataforma que recebe seu código-fonte e o executa \\
\textbf{SaaS} & só os \textbf{dados} e a configuração de uso & software
pronto, acessado pelo navegador \\
\bottomrule
\end{tabular}
\end{center}

Duas consequências importantes. Primeira: \textbf{quanto mais “as a
Service”}, \textbf{menos controle} você tem --- não é só menos trabalho
operacional, é também menos liberdade para customizar a infraestrutura
subjacente. Segunda, a \textbf{responsabilidade compartilhada em
segurança}: mesmo em SaaS, \textbf{o dado continua sendo seu} --- o
provedor responde pela segurança \textit{da} nuvem (a infraestrutura em
si), e o cliente responde pela segurança \textit{na} nuvem (quem tem
acesso, como o dado é classificado, se a conta usa segundo fator de
autenticação). A afirmação de que SaaS “transfere integralmente a
responsabilidade pelos dados ao provedor” é falsa.
\end{conceito}

\begin{conceito}[Escalabilidade horizontal x vertical]
\textbf{Escalabilidade vertical (\textit{scale up})}: adicionar recursos à
\textbf{instância existente} (mais CPU, mais RAM na mesma máquina) ---
tem um teto físico e costuma exigir reinício.
\textbf{Escalabilidade horizontal (\textit{scale out})}: \textbf{adicionar
mais instâncias/nós} rodando em paralelo --- teoricamente sem teto, mas
exige que a aplicação seja capaz de distribuir carga entre instâncias (o
que REST \textit{stateless} viabiliza diretamente).

\textbf{Elasticidade}: ajustar a capacidade \textbf{automaticamente}
conforme a demanda sobe e desce, sem intervenção manual. \textbf{Cloud
bursting}: uma aplicação que roda normalmente em infraestrutura própria
“estoura” para a nuvem pública só nos picos de demanda, voltando ao normal
depois.
\end{conceito}

\begin{conceito}[Serverless / FaaS]
“Sem servidor” é um nome enganoso: \textbf{existe} servidor, só que é o
\textbf{provedor} quem o gerencia, não o cliente. O código é publicado como
uma \textbf{função} e o provedor a executa em resposta a \textbf{eventos}
(uma requisição HTTP, uma mensagem numa fila, um arquivo novo num
\textit{bucket}).

Características cobradas em prova: escala \textbf{de zero} até muitas
execuções simultâneas conforme os eventos chegam; \textbf{cobrança pelo
tempo e pelos recursos efetivamente consumidos} durante a execução (não
por instância ficando parada); execução \textbf{sem estado} entre uma
invocação e a próxima (o estado precisa ir para um armazenamento externo,
como um banco); \textbf{\textit{cold start}} --- a latência maior na
primeira invocação depois de um período ocioso, porque o ambiente de
execução precisa ser inicializado; e um \textbf{limite de tempo por
execução}, que desaconselha serverless para processamento longo e
contínuo.
\end{conceito}

\begin{exemplo}[Balanceador de carga x CDN --- gargalos diferentes]
Um site de e-commerce está lento para \textbf{todos} os usuários porque as
poucas instâncias da aplicação estão sobrecarregadas de processamento
$\to$ a solução é um \textbf{balanceador de carga}, que reparte as
requisições entre mais instâncias. O mesmo site está rápido para usuários
no Brasil, mas lento para usuários na Europa, ao carregar imagens e
arquivos estáticos $\to$ a solução é uma \textbf{CDN}, que replica esse
conteúdo estático em pontos de presença geograficamente distribuídos e o
entrega do nó mais próximo do usuário. Os dois atacam gargalos diferentes
(processamento x latência geográfica de conteúdo estático) e por isso se
somam na prática, em vez de serem alternativas mutuamente exclusivas.
\end{exemplo}

\begin{cuidado}
“Adicionar recursos à instância” é sempre \textbf{vertical}, nunca
horizontal. Em serverless: “não há servidor, o código roda no dispositivo
que dispara o evento” é falso (há servidor, do lado do provedor); “o
estado da execução anterior permanece disponível” é falso (é sem estado
entre invocações); “a latência da primeira invocação é igual à das
demais” ignora o \textit{cold start}; “remove o limite de tempo por
execução” é falso --- o limite existe e é justamente o que desaconselha
processamento longo em FaaS.
\end{cuidado}

\subsection{Containers e virtualização}

\begin{conceito}
\textbf{Container (Docker) x máquina virtual}: container compartilha o
\textbf{kernel} do sistema operacional hospedeiro (por isso é leve e sobe
rápido); máquina virtual tem um \textbf{sistema operacional convidado
completo}, rodando sobre um \textbf{hipervisor} (isola mais, mas pesa mais
e sobe mais devagar). \textbf{Kubernetes} orquestra containers em escala
(agenda onde cada container roda, reinicia os que caem, distribui carga
entre eles).

\textbf{Hipervisor Tipo 1} (\textit{bare-metal}, roda diretamente sobre o
hardware, sem sistema operacional hospedeiro no meio: ESXi, Hyper-V, KVM)
x \textbf{Tipo 2} (\textit{hosted}, roda como um programa em cima de um
sistema operacional já instalado: VirtualBox, VMware Workstation).

\textbf{Virtualização total} (o sistema operacional convidado não sabe que
está virtualizado, roda sem modificação) x \textbf{paravirtualização} (o
sistema operacional convidado é modificado para cooperar diretamente com o
hipervisor via \textit{hypercalls}, o que é mais rápido por evitar
simulação completa de hardware).
\end{conceito}

\begin{cuidado}
Container \textbf{não} é sinônimo de VM --- a diferença central é
compartilhar o kernel (container) versus ter sistema operacional próprio
completo (VM).
\end{cuidado}

\section{Transações distribuídas}
\label{sec:transacoes-dist}

\begin{conceito}[O problema nasce quando cada serviço ganha o próprio banco]
Dentro de um único SGBD, o ACID (Capítulo \ref{cap:banco-dados-teoria}) é
“de graça”: existe um só log de transações, um só gerenciador de
bloqueios, e ele decide sozinho se a transação efetiva ou desfaz.
Distribuído entre vários bancos (um por microsserviço), não há árbitro
único --- “debitar no serviço de contas e reservar no serviço de estoque”
são \textbf{duas} transações, em \textbf{dois} bancos diferentes, ligadas
por uma rede que pode falhar no meio do caminho. Não existe uma operação
que efetive as duas ao mesmo tempo de forma atômica; o que existe são
\textbf{protocolos que tentam coordenar} isso, cada um pagando um preço
diferente.

\begin{itemize}
\item \textbf{2PC (\textit{Two-Phase Commit})}: um coordenador central
pergunta a todos os participantes se estão prontos para efetivar; só faz o
\textit{commit} se \textbf{todos} confirmarem (unanimidade). É
\textbf{bloqueante}: se o coordenador cair no meio do processo, os
participantes ficam travados esperando, sem saber se devem efetivar ou
desfazer.
\item \textbf{Saga}: uma sequência de transações locais independentes, cada
uma com uma ação de \textbf{compensação} correspondente caso algo dê
errado mais adiante na sequência (se reservar estoque funcionou mas debitar
o cartão falhou, uma ação compensatória libera o estoque reservado). É o
padrão típico para transações que atravessam microsserviços.
\item \textbf{Raft}: algoritmo de \textbf{consenso}, que resolve um
problema \textbf{diferente} do 2PC --- não pergunta “todos aceitam
efetivar esta transação?”, pergunta “em que \textbf{sequência de
operações} este grupo de réplicas concorda?”. Funciona por
\textbf{eleição de líder} (o líder ordena as operações e replica o log
para os demais) e decide por \textbf{maioria} (quórum) --- o que o torna
\textbf{tolerante a falhas}: o sistema sobrevive à queda de uma minoria
das réplicas, inclusive do próprio líder, que é substituído por uma nova
eleição. É o algoritmo que sustenta o \textit{etcd} (e, por consequência,
o próprio Kubernetes).
\end{itemize}
\end{conceito}

\begin{cuidado}[2PC x Raft --- o par central desta seção]
O 2PC é \textbf{commit atômico}: exige \textbf{unanimidade} e
\textbf{bloqueia} se o coordenador cair (ponto único de falha). O Raft é
\textbf{consenso}: decide por \textbf{maioria} e \textbf{continua
funcionando} mesmo com uma minoria de réplicas fora do ar. Trocar
“unanimidade” por “maioria” de um algoritmo para o outro é o erro mais
comum ao descrevê-los, e afirmar que o 2PC “tolera a falha do coordenador”
é falso --- é exatamente a fraqueza estrutural do protocolo.
\end{cuidado}

\section{Tópicos complementares}

\begin{conceito}[DDD --- Domain-Driven Design]
Abordagem de modelagem que organiza o software em torno do
\textbf{domínio de negócio}, não da tecnologia. Conceitos centrais:
\textbf{Aggregate} (fronteira de consistência transacional --- protege as
invariantes de um conjunto de entidades relacionadas, e não expõe
referência direta a cada entidade interna a quem está fora do agregado),
\textbf{Repository} (abstração de coleção/persistência do agregado, sem
que classes externas instanciem entidades internas diretamente),
\textbf{Factory} (encapsula a criação de objetos complexos do domínio),
\textbf{Ubiquitous Language} (vocabulário compartilhado entre
especialistas de domínio e código-fonte --- o nome da classe no código
deve ser o mesmo termo que o especialista de negócio usa). Eventos de
domínio são tipicamente \textbf{imutáveis}, porque registram um fato que
já ocorreu no passado (não faz sentido “editar” um fato passado).
\end{conceito}

\begin{conceito}[Service mesh x API gateway]
\textbf{API gateway}: ponto único de entrada dos microsserviços,
cuidando do tráfego \textbf{norte-sul} (de fora para dentro do sistema) ---
roteamento, autenticação, limitação de taxa, agregação de respostas.
\textbf{Service mesh} (ex.: Istio): cuida da comunicação
\textbf{serviço-a-serviço}, tráfego \textbf{leste-oeste} (dentro do
sistema, entre os próprios microsserviços) --- descoberta de serviço,
balanceamento, retentativa, criptografia entre serviços, tudo de forma
transparente à aplicação.
\end{conceito}

\begin{conceito}[Kubernetes: taint e toleration]
Funcionam de um jeito contraintuitivo em relação ao nome. O
\textbf{taint} é aplicado ao \textbf{nó} e \textbf{afasta} pods: “não
escalone nada aqui, a menos que aceite esta marca”. A \textbf{toleration}
é aplicada ao \textbf{pod} e diz “eu tolero essa marca”, tornando o pod
\textbf{elegível} para aquele nó --- útil para reservar nós especiais
(com GPU, banco de dados, ou o nó mestre) só para cargas de trabalho
específicas. Tolerar \textbf{não é atrair}: um pod que tolera o taint
\textbf{pode} ser escalonado para lá, não é \textbf{obrigado} a ir --- quem
efetivamente atrai um pod para um nó específico é \textit{node
affinity}/\textit{nodeSelector}, um mecanismo diferente.
\end{conceito}

\begin{conceito}[IaC --- Infraestrutura como Código]
Provisiona ambiente de infraestrutura de forma \textbf{declarativa}
(descreve o estado desejado, não o passo a passo) e \textbf{versionada}
(o ambiente vira arquivo de texto sob controle de versão, revisável e
reproduzível). Ferramentas típicas: Terraform, Ansible.
\end{conceito}
